% =============================================================================
% Chapter 2: Related Work (Background) -- Part I
% =============================================================================
\mychapter{2. Related Work (Background)}

This chapter presents the theoretical foundations and related work underpinning the development of
an automated defect detection pipeline for XCT data of additively manufactured metal components. It
covers the physics of defect formation in AM processes, the principles of XCT imaging and its
associated artefacts, classical and deep learning-based detection methods, and the mathematical
foundations of the loss functions and evaluation metrics used in this thesis.

\mysection{2.1 Additive Manufacturing and Internal Defect Formation}
Metal additive manufacturing encompasses a family of layer-by-layer fabrication processes in which
components are built from metallic powder or wire feedstock using a focused energy
source~\citep{Smith2022}. The most industrially prevalent process is laser powder bed fusion (LPBF),
in which a laser selectively melts successive layers of powder spread across a build platform. The
rapid solidification dynamics, characterised by cooling rates on the order of $10^{6}$~K/s, create
complex microstructures and residual stress distributions that differ substantially from those
produced by conventional casting~\citep{Zhang2021}.

Internal defects arise from several distinct physical mechanisms. Gas porosity forms when dissolved
gas is entrapped in the melt pool; these defects are typically spherical with diameters ranging from
a few micrometres to several hundred micrometres~\citep{Zhang2021}. Lack-of-fusion (LoF) defects
arise when adjacent scan tracks are incompletely fused due to insufficient energy density; they are
irregularly shaped with sharp features that act as stress concentrators~\citep{Smith2022}.
Micro-cracks can form due to solidification cracking or solid-state cracking driven by thermal
residual stresses. Inclusions may arise from surface oxide films on powder particles, redeposited
condensate, or contamination~\citep{Deshpande2024}.

The mechanical consequences of these defects are severe. Even small pores or LoF voids can initiate
fatigue crack growth under cyclic loading, dramatically reducing component service life. In
safety-critical applications such as aerospace turbine blades or biomedical implants, undetected
defects can lead to catastrophic failure~\citep{Zubayer2023}. This motivates the need for reliable,
automated, and scalable defect detection methods.

\mysection{2.2 X-ray Computed Tomography for Non-Destructive Inspection}
X-ray computed tomography is a non-destructive imaging technique that reconstructs the
three-dimensional internal structure of an object from a series of two-dimensional X-ray projection
images acquired at multiple angles~\citep{Miller2020}. The achievable voxel size in an XCT system is
governed by the geometric magnification relationship:
\begin{equation}
V_{xl} = \frac{\mathrm{SOD}}{\mathrm{SDD}} \times P_{xl}
\label{eq:voxel_size}
\end{equation}
where SOD is the source-to-object distance, SDD is the source-to-detector distance, and $P_{xl}$ is
the detector pixel pitch. For metallic AM components, typical XCT voxel sizes range from
1--50~$\mu$m~\citep{Miller2020}.

The NIST CoCr AM XCT dataset used in this thesis provides a representative example of
industrial-scale XCT data: 900 slices at $1010 \times 980$ pixel resolution, stored as 16-bit
unsigned TIFF images with a raw intensity range of $[0, 59631]$~\citep{NIST_CoCr}. At this scale,
manual inspection of all slices is impractical, motivating automated detection approaches.

\mysubsection{2.2.1 Imaging Artefacts in Metal XCT}
Several systematic imaging artefacts complicate automated defect detection in dense metal XCT
scans~\citep{Miller2020}:
\begin{itemize}
\item \textbf{Beam hardening}: Differential attenuation of low-energy X-ray photons creates a
  cupping effect, a systematic intensity gradient where the component periphery appears brighter
  than the centre. This gradient can mask real defects near component boundaries or create false
  intensity variations that mimic defects.
\item \textbf{Ring artefacts}: Non-uniformities in individual detector elements produce concentric
  ring patterns in reconstructed slices. These rings interfere with defect localisation by creating
  circular false-positive regions at fixed radial distances from the rotation centre.
\item \textbf{Noise}: Quantum noise from photon counting statistics and electronic detector noise
  introduce random intensity fluctuations that reduce contrast between defect voids and the
  surrounding metal matrix.
\end{itemize}
Addressing these artefacts through preprocessing is a prerequisite for reliable automated defect
detection, as classical and deep learning methods alike are sensitive to systematic intensity
variations~\citep{Bellens2021}.

\mysection{2.3 Preprocessing Methods for XCT Data}
\mysubsection{2.3.1 Percentile Normalisation}
Raw XCT intensity values vary widely across scans and scanners due to differences in energy settings,
detector gain, and material composition. Percentile normalisation clips intensities to a robust range
and rescales to $[0, 1]$:
\begin{equation}
x_{\mathrm{norm}} = \frac{\mathrm{clip}(x, p_1, p_{99}) - p_1}{p_{99} - p_1 + \varepsilon}
\label{eq:norm}
\end{equation}
where $p_1$ and $p_{99}$ are the 1st and 99th intensity percentiles and $\varepsilon = 10^{-7}$
prevents division by zero. This approach is robust to extreme outliers such as metal streak artefacts
and hot pixels that would distort linear normalisation~\citep{Bellens2021}.

\mysubsection{2.3.2 Beam Hardening Correction}
Beam hardening correction estimates the systematic intensity trend along the beam axis by fitting a
polynomial of degree $d$ to the mean intensity profile and subtracting it~\citep{Miller2020}:
\begin{equation}
x_{\mathrm{BHC}}[i] = x_{\mathrm{norm}}[i] - \big(T(i) - \bar{T}\big),
\qquad T(i) = \sum_{k=0}^{d} a_k\, i^{k}
\label{eq:bhc}
\end{equation}
where $T(i)$ is the fitted polynomial trend and $\bar{T}$ is its mean, ensuring zero net intensity
shift. A polynomial degree of $d = 3$ was selected through parameter sensitivity analysis on the NIST
CoCr dataset, as lower degrees are insufficient to capture the non-linear cupping gradient while
higher degrees risk over-fitting to slice-level intensity variations.

\mysubsection{2.3.3 Ring Artefact Suppression}
Ring artefacts appear as horizontal bands in the polar coordinate representation of each slice. The
suppression algorithm converts each slice to polar coordinates, applies a median filter of radius $r$
pixels along the angular axis, and converts back to Cartesian coordinates. A filter radius of
$r = 15$ pixels was found optimal for the NIST CoCr dataset, achieving effective ring suppression
without introducing spatial blurring that would obscure small defects.

\mysubsection{2.3.4 Non-Local Means Denoising}
Non-local means (NLM) denoising~\citep{Buades2005} replaces each voxel intensity with a weighted
average of voxels in a search window, where weights are determined by patch-level intensity
similarity:
\begin{equation}
\mathrm{NLM}(x)_i = \sum_j w(i,j)\, x_j,
\qquad
w(i,j) = \frac{\exp\!\big(-\lVert P_i - P_j \rVert^2 / h^2\big)}{Z_i}
\label{eq:nlm}
\end{equation}
where $P_i$ and $P_j$ are intensity patches centred at voxels $i$ and $j$, $h$ is the filter strength
parameter, and $Z_i$ is a normalisation constant. Unlike Gaussian or median filters, NLM preserves
sharp defect boundaries because it averages over structurally similar regions rather than spatially
adjacent ones~\citep{Buades2005}. Applied to the NIST CoCr dataset with $h = 0.6 \times \hat{\sigma}$
(adaptive per slice; corrected from the mid-term submission's stated $1.0 \times \hat{\sigma}$, see
the Note on this combined document), NLM reduces mean noise standard deviation from 20952 to 0.385, a
factor of $\times 54{,}500$.

\mysection{2.4 Classical Defect Detection Methods}
\mysubsection{2.4.1 Otsu Thresholding}
The Otsu method~\citep{Otsu1979} is the most widely used classical approach for XCT defect detection.
Given a normalised grey-level histogram $h(i) = n_i / n$, the optimal threshold $t^*$ maximises the
inter-class variance:
\begin{equation}
\sigma_B^2(t) = \omega_1(t)\, \omega_2(t)\, \big[\mu_2(t) - \mu_1(t)\big]^2,
\qquad t^* = \arg\max_t \sigma_B^2(t)
\label{eq:otsu}
\end{equation}
where $\omega_1(t)$ and $\omega_2(t)$ are the class probabilities below and above threshold $t$, and
$\mu_1(t), \mu_2(t)$ are the corresponding class means~\citep{Otsu1979}.

Otsu thresholding is fully automatic and computationally efficient, but fails under conditions
typical of metal XCT: beam hardening creates systematic intensity gradients that shift the optimal
threshold spatially, and low contrast between defect voids and the metal matrix limits
separability~\citep{Bellens2021}. On the NIST AM Bench dataset, Otsu achieves Dice scores of only
0.43--0.61 for porosity detection, compared to 0.82--0.91 for CNN-based approaches~\citep{Wong2022}.
In this thesis, Otsu is used both as a classical detection baseline and as an automatic label
generator for U-Net training supervision.

\mysubsection{2.4.2 Morphological Post-Processing}
Raw Otsu masks are refined through morphological operations to improve label
quality~\citep{Deshpande2024}:
\begin{itemize}
\item Morphological opening ($3\times3$ structuring element): removes isolated noise voxels that
  survive thresholding
\item Morphological closing ($3\times3$ structuring element): fills small holes within genuine
  defect regions
\item Connected component filtering: discards components smaller than 5 voxels as noise artefacts
\end{itemize}

\mysection{2.5 Deep Learning for Defect Detection}
\mysubsection{2.5.1 Convolutional Neural Networks}
Deep convolutional neural networks (CNNs) have demonstrated transformative performance across image
analysis tasks due to their translational equivariance and hierarchical feature
learning~\citep{Smith2022}. A convolutional layer applies $K$ learnable filters of spatial extent $F$
to produce output feature maps. The output dimensions are:
\begin{equation}
H_2 = \frac{H_1 - F + 2P}{S} + 1,
\qquad
W_2 = \frac{W_1 - F + 2P}{S} + 1
\label{eq:conv_dims}
\end{equation}
where $H_1, W_1$ are input spatial dimensions, $P$ is zero padding, and $S$ is stride. The ReLU
activation function introduces non-linearity at low computational cost:
\begin{equation}
\mathrm{ReLU}(x) = \max(0, x)
\label{eq:relu}
\end{equation}

\mysubsection{2.5.2 U-Net Architecture}
The U-Net architecture~\citep{Ronneberger2015} was originally developed for biomedical image
segmentation and has since become the dominant approach for defect detection in XCT
data~\citep{Deshpande2024}. It consists of a symmetric encoder-decoder structure with skip
connections that transfer spatial detail from encoder to decoder:
\begin{itemize}
\item \textbf{Encoder}: four downsampling stages, each comprising two $3\times3$ convolutional
  layers with batch normalisation and ReLU activation~\citep{Ronneberger2015}, followed by
  $2\times2$ max pooling. Feature channels double at each stage: $64 \to 128 \to 256 \to 512$.
\item \textbf{Bottleneck}: two $3\times3$ convolutional layers with 1024 channels and dropout
  regularisation ($p = 0.2$) to reduce overfitting on limited training data.
\item \textbf{Decoder}: four upsampling stages with transposed convolutions and skip connection
  concatenation from the corresponding encoder level.
\item \textbf{Output layer}: $1\times1$ convolution followed by sigmoid activation, producing
  per-voxel defect probability in $[0, 1]$:
\begin{equation}
\sigma(x) = \frac{1}{1 + e^{-x}}
\label{eq:sigmoid}
\end{equation}
\end{itemize}
Skip connections are critical for detecting small defects, as they preserve high-resolution spatial
detail that is lost during downsampling~\citep{Ronneberger2015}. The 2D formulation processes
individual XCT slices as $256 \times 256$ pixel patches, achieving within 2--4\% of full 3D
volumetric performance at 10--50$\times$ lower computational cost~\citep{Wong2021}.

\mysubsection{2.5.3 Automatic Label Generation}
A critical challenge in this thesis is the absence of expert-annotated ground truth masks for the
NIST CoCr dataset. This is addressed through a two-stage automatic label generation strategy
following Yosifov~\citep{Yosifov2020}:
\begin{enumerate}
\item A first-pass Otsu threshold separates the metal component from the air background, defining
  the region of interest
\item A second-pass Otsu threshold is applied within the metal region only, identifying dark voxels
  as defects
\end{enumerate}
This two-stage approach is essential to avoid misclassifying the air background outside the
cylindrical component as defects, a common failure mode of single-pass global thresholding on XCT
data~\citep{Otsu1979}. The resulting pseudo-labels are imperfect but sufficient to provide training
supervision, with the U-Net learning to generalise beyond the label generator's systematic
errors~\citep{Yosifov2020}.

\mysection{2.6 Loss Functions for Class-Imbalanced Detection}
A central challenge in XCT defect detection is severe class imbalance: defect voxels constitute only
1.2--2.8\% of total scan volume~\citep{Sudre2017}. Standard binary cross-entropy (BCE) loss causes
models to collapse to predicting background for all voxels, achieving high accuracy while detecting
no defects:
\begin{equation}
\mathcal{L}_{\mathrm{BCE}} = -\frac{1}{N} \sum_{i=1}^{N}
\Big[ g_i \log(p_i) + (1 - g_i) \log(1 - p_i) \Big]
\label{eq:bce}
\end{equation}
where $p_i$ is the predicted defect probability and $g_i$ is the ground truth label for voxel $i$.

\mysubsection{2.6.1 Dice Loss}
Dice loss directly optimises the overlap between predicted and true defect regions, making it
inherently robust to class imbalance~\citep{Sudre2017}:
\begin{equation}
\mathcal{L}_{\mathrm{Dice}} = 1 -
\frac{2 \sum_i p_i g_i + \varepsilon}{\sum_i p_i + \sum_i g_i + \varepsilon}
\label{eq:dice_loss}
\end{equation}
where $\varepsilon$ is a smoothing constant that prevents division by zero.

\mysubsection{2.6.2 Focal Loss}
Focal loss~\citep{Lin2017} addresses class imbalance by down-weighting easy background examples and
focusing training on hard defect voxels:
\begin{equation}
\mathcal{L}_{\mathrm{Focal}} = -\alpha (1 - p_t)^{\gamma} \log(p_t),
\qquad \alpha = 0.25,\ \gamma = 2
\label{eq:focal_loss}
\end{equation}
where $(1 - p_t)^{\gamma}$ is the modulating factor that reduces the loss contribution of confidently
classified background voxels.

\mysubsection{2.6.3 Combined Dice-Focal Loss}
The primary loss formulation used in this thesis combines Dice and Focal losses with equal
weighting~\citep{Sudre2017,Lin2017}:
\begin{equation}
\mathcal{L}_{\mathrm{DF}} = \lambda \mathcal{L}_{\mathrm{Dice}} + (1 - \lambda) \mathcal{L}_{\mathrm{Focal}},
\qquad \lambda = 0.5
\label{eq:dice_focal}
\end{equation}
This formulation benefits from Dice loss's class-imbalance robustness and Focal loss's hard-example
mining, achieving a validation Dice score of 0.68 after 20 epochs compared to 0.41 for BCE
alone~\citep{Sudre2017}.

\mysection{2.7 Evaluation Metrics}
Detection performance is quantified using the following standard metrics:
\begin{align}
\mathrm{DSC} &= \frac{2 |A \cap B|}{|A| + |B|} \label{eq:dsc} \\[4pt]
\mathrm{IoU} &= \frac{|A \cap B|}{|A \cup B|} \label{eq:iou} \\[4pt]
\mathrm{Precision} &= \frac{TP}{TP + FP}, \qquad
\mathrm{Recall} = \frac{TP}{TP + FN} \label{eq:prec_rec}
\end{align}
where $A$ is the predicted detection mask, $B$ is the reference label, $TP$ is true positives, $FP$
is false positives, and $FN$ is false negatives. The Dice coefficient~\citep{Otsu1979} is the primary
acceptance criterion, with a target of $\mathrm{DSC} \geq 0.75$ on the held-out test set.

\mysection{2.8 Data Augmentation}
Data augmentation is essential for training robust defect detection models when labelled data is
limited~\citep{Bellens2021}. By applying random transformations to training patches, augmentation
effectively increases dataset diversity and reduces overfitting to scan-specific characteristics. The
augmentation operations applied in this thesis are summarised in Table~\ref{tab:augmentation_ops}.

\begin{table}[htbp]
\centering
\caption{Data augmentation operations applied during training~\citep{Bellens2021}.}
\label{tab:augmentation_ops}
\begin{tabular}{lll}
\toprule
\textbf{Augmentation} & \textbf{Probability} & \textbf{Parameters} \\
\midrule
Horizontal/vertical flip & 0.5 & --- \\
Random $90^\circ$ rotation & 0.5 & $k \in \{1, 2, 3\}$ \\
Elastic deformation & 0.3 & $\alpha = 34,\ \sigma = 4$ \\
Intensity scaling & 0.5 & Scale $\in [0.9, 1.1]$ \\
Gaussian noise & 0.5 & $\sigma \in [0.01, 0.05]$ \\
Gamma correction & 0.5 & $\gamma \in [0.8, 1.2]$ \\
\bottomrule
\end{tabular}
\end{table}

All transforms are applied identically to both image patch and corresponding label patch to maintain
spatial alignment. Elastic deformation simulates the geometric variability of defect shapes across
different scan positions, while intensity transforms simulate acquisition variability across
different scan settings and materials~\citep{Bellens2021}.

\mysection{2.9 Related Work Summary}
Table~\ref{tab:related_work_summary} summarises the key findings from the literature that directly
inform the design decisions in this thesis.

\begin{table}[htbp]
\centering
\caption{Summary of directly relevant related work.}
\label{tab:related_work_summary}
\small
\begin{tabular}{p{2.6cm}p{5.4cm}p{5.4cm}}
\toprule
\textbf{Study} & \textbf{Finding} & \textbf{Relevance to this thesis} \\
\midrule
Wong et al.~\citep{Wong2022} & CNN achieves Dice 0.82--0.91 vs.\ Otsu 0.43--0.61 & Motivates U-Net
  over classical baseline \\
Bellens et al.~\citep{Bellens2021} & Augmentation essential for limited labelled data & Justifies
  six-transform augmentation pipeline \\
Zubayer et al.~\citep{Zubayer2023} & Data diversity $>$ architecture complexity & Supports 2D-first
  approach on diverse patches \\
Liu et al.~\citep{Liu2024} & Label quality is dominant performance factor & Justifies morphological
  cleaning of Otsu labels \\
Deshpande et al.~\citep{Deshpande2024} & Standardised benchmarks are the primary gap & Motivates use
  of NIST CoCr as benchmark dataset \\
Yosifov~\citep{Yosifov2020} & Otsu pseudo-labels enable U-Net training without annotation & Basis
  for automatic label generation strategy \\
Wong et al.~\citep{Wong2021} & 2D within 2--4\% of 3D at 10--50$\times$ lower cost & Justifies 2D
  slice-based detection approach \\
Sudre et al.~\citep{Sudre2017} & Combined Dice-Focal loss best for class imbalance & Primary loss
  function selection \\
\bottomrule
\end{tabular}
\end{table}
